\sectionlink*{索引}
\begingroup\small
\setlength{\parindent}{0pt}
\setlength{\parskip}{1pt}
\newcommand{\idxentry}[2]{\par\noindent\hangindent=1em\hangafter=1\textbf{#1}\quad#2\par}
\newcommand{\idxsub}[2]{\par\noindent\hangindent=2em\hangafter=1\hspace*{1em}#1\quad#2\par}
\begin{multicols}{2}
\raggedcolumns
\idxentry{観測による現実化}{\sourcepageref{19}}
\idxentry{補助系（アンシラ）}{\sourcepageref{311}--\sourcepageref{320}}
\idxentry{異常拡散係数}{\sourcepageref{186}, \sourcepageref{191}, \sourcepageref{200}, \sourcepageref{205}, \sourcepageref{220}}
\idxentry{装置}{\sourcepageref{32}, \sourcepageref{49}, \sourcepageref{51}, \sourcepageref{320}--\sourcepageref{321}}
\idxentry{人工神経回路網}{\sourcepageref{365}}
\idxentry{非対称エネルギー}{\sourcepageref{174}}
\idxentry{軸索}{\sourcepageref{368}}
\idxentry{浴}{\sourcepageref{171}}
\idxentry{ビアブル（beables）}{\sourcepageref{16}}
\idxentry{Bell不等式}{\sourcepageref{17}}
\idxentry{Bell状態}{\sourcepageref{29}, \sourcepageref{32}}
\idxentry{双線形結合}{\sourcepageref{179}}
\idxentry{ビット反転誤り}{\sourcepageref{309}}
\idxentry{測定のブラックボックス的見方}{\sourcepageref{334}}
\idxentry{Bloch球}{\sourcepageref{295}}
\idxentry{Bohm力学}{\sourcepageref{18}, \sourcepageref{354}--\sourcepageref{357}}
\idxsub{デコヒーレンスとの関係}{\sourcepageref{356}--\sourcepageref{357}}
\idxsub{粒子概念}{\sourcepageref{356}}
\idxsub{軌道}{\sourcepageref{355}}
\idxentry{Born近似}{\sourcepageref{156}, \sourcepageref{160}--\sourcepageref{161}}
\idxentry{Born確率}{\sourcepageref{35}, \sourcepageref{333}}
\idxentry{Born則}{\sourcepageref{15}, \sourcepageref{35}, \sourcepageref{331}}
\idxsub{エンバリアンスとの関係}{\sourcepageref{341}--\sourcepageref{343}}
\idxsub{Zurekによる導出}{\sourcepageref{341}--\sourcepageref{343}}
\idxentry{Born–Markovマスター方程式}{\sourcepageref{155}--\sourcepageref{165}}
\idxsub{導入される仮定}{\sourcepageref{156}}
\idxsub{導出}{\sourcepageref{158}--\sourcepageref{165}}
\idxsub{構造}{\sourcepageref{156}--\sourcepageref{158}}
\idxentry{Born–Pauli解釈}{\sourcepageref{19}}
\idxentry{Bose–Einstein凝縮体}{\sourcepageref{282}--\sourcepageref{283}}
\idxentry{脳}{\sourcepageref{359}--\sourcepageref{378}}
\idxentry{分枝}{\sourcepageref{337}, \sourcepageref{375}}
\idxentry{$\mathrm C_{70}$の干渉実験}{\sourcepageref{258}--\sourcepageref{270}}
\idxsub{熱放射の放出}{\sourcepageref{265}--\sourcepageref{267}}
\idxsub{干渉模様}{\sourcepageref{262}--\sourcepageref{263}}
\idxsub{実験配置}{\sourcepageref{259}--\sourcepageref{262}}
\idxsub{経路情報}{\sourcepageref{263}--\sourcepageref{265}}
\idxentry{Caldeira–Leggettマスター方程式}{\sourcepageref{191}--\sourcepageref{194}}
\idxentry{標準模型}{\sourcepageref{171}--\sourcepageref{241}}
\idxentry{共振器量子電気力学}{\sourcepageref{255}}
\idxentry{電荷量子ビット}{\sourcepageref{281}}
\idxentry{キラル分子}{\sourcepageref{83}}
\idxentry{古典的集団}{\sourcepageref{37}}
\idxentry{CNOTゲート}{\sourcepageref{306}}
\idxentry{粗視化}{\sourcepageref{37}, \sourcepageref{59}}
\idxentry{符号部分空間}{\sourcepageref{315}}
\idxentry{コヒーレンス}{\sourcepageref{5}, \sourcepageref{20}}
\idxsub{非局在化}{\sourcepageref{8}}
\idxsub{再帰}{\sourcepageref{92}, \sourcepageref{213}}
\idxsub{再局在化}{\sourcepageref{99}, \sourcepageref{255}--\sourcepageref{258}}
\idxsub{復活}{\sourcepageref{98}--\sourcepageref{101}, \sourcepageref{255}--\sourcepageref{258}}
\idxentry{コヒーレンス長}{\sourcepageref{140}--\sourcepageref{147}}
\idxentry{コヒーレント振動}{\sourcepageref{275}--\sourcepageref{279}}
\idxentry{コヒーレントな広がり}{\sourcepageref{116}, \sourcepageref{144}, \sourcepageref{145}}
\idxentry{コヒーレント状態}{\sourcepageref{117}, \sourcepageref{202}, \sourcepageref{249}}
\idxentry{収縮公準}{\sourcepageref{3}, \sourcepageref{15}, \sourcepageref{331}}
\idxentry{可換性基準}{\sourcepageref{77}}
\idxentry{相補性原理}{\sourcepageref{61}--\sourcepageref{65}}
\idxentry{条件付き量子動力学}{\sourcepageref{306}}
\idxentry{共役基底}{\sourcepageref{23}}
\idxentry{意識}{\sourcepageref{363}--\sourcepageref{367}, \sourcepageref{370}--\sourcepageref{371}, \sourcepageref{376}--\sourcepageref{378}}
\idxentry{保存則}{\sourcepageref{30}}
\idxentry{文脈依存性}{\sourcepageref{18}}
\idxentry{連続的自発的局在化}{\sourcepageref{349}--\sourcepageref{354}}
\idxentry{量子コンピューターにおける制御可能性}{\sourcepageref{301}--\sourcepageref{302}}
\idxentry{Cooper対}{\sourcepageref{271}}
\idxentry{Copenhagen解釈}{\sourcepageref{330}--\sourcepageref{336}, \sourcepageref{360}, \sourcepageref{363}}
\idxentry{相関}{}
\idxsub{古典的}{\sourcepageref{30}}
\idxsub{量子的}{\sourcepageref{29}}
\idxsub{量子的相関と古典的相関}{\sourcepageref{30}--\sourcepageref{32}}
\idxentry{dc-SQUID}{\sourcepageref{272}}
\idxentry{de Broglie波長（物質）}{\sourcepageref{116}}
\idxentry{de Broglie波長（熱的）}{\sourcepageref{94}}
\idxentry{デコヒーレンス}{}
\idxsub{量子誤り訂正との関係}{\sourcepageref{311}--\sourcepageref{321}}
\idxsub{部分系への分割との関係}{\sourcepageref{101}--\sourcepageref{103}}
\idxsub{装置が誘起するもの}{\sourcepageref{320}--\sourcepageref{321}}
\idxsub{標準模型}{\sourcepageref{171}--\sourcepageref{241}}
\idxsub{集団的}{\sourcepageref{322}--\sourcepageref{323}}
\idxsub{環境散乱によるもの}{\sourcepageref{115}--\sourcepageref{151}}
\idxsub{散乱誘起デコヒーレンスの動力学}{\sourcepageref{139}--\sourcepageref{150}}
\idxsub{実験}{\sourcepageref{243}--\sourcepageref{291}}
\idxsub{因子}{\sourcepageref{91}, \sourcepageref{122}, \sourcepageref{212}}
\idxsub{偽のもの}{\sourcepageref{96}}
\idxsub{歴史}{\sourcepageref{10}--\sourcepageref{12}}
\idxsub{Bose–Einstein凝縮体におけるもの}{\sourcepageref{282}--\sourcepageref{283}}
\idxsub{共振器QEDにおけるもの}{\sourcepageref{251}--\sourcepageref{258}}
\idxsub{エネルギーでのもの}{\sourcepageref{81}, \sourcepageref{228}}
\idxsub{量子力学の解釈におけるもの}{\sourcepageref{329}--\sourcepageref{358}}
\idxsub{運動量でのもの}{\sourcepageref{194}}
\idxsub{ニューロンにおけるもの}{\sourcepageref{368}--\sourcepageref{371}}
\idxsub{位相空間でのもの}{\sourcepageref{173}, \sourcepageref{194}--\sourcepageref{203}}
\idxsub{量子Brown運動におけるもの}{\sourcepageref{178}--\sourcepageref{206}}
\idxsub{量子コンピューターにおけるもの}{\sourcepageref{293}--\sourcepageref{328}}
\idxsub{量子電気機械系におけるもの}{\sourcepageref{287}--\sourcepageref{288}}
\idxsub{SQUIDにおけるもの}{\sourcepageref{279}--\sourcepageref{282}}
\idxsub{脳におけるもの}{\sourcepageref{365}--\sourcepageref{374}}
\idxsub{独立なもの}{\sourcepageref{322}--\sourcepageref{323}}
\idxsub{不可逆性}{\sourcepageref{5}, \sourcepageref{69}}
\idxsub{動力学的時間尺度でのもの}{\sourcepageref{197}}
\idxsub{率}{$\to$ デコヒーレンス率}
\idxsub{逆転}{\sourcepageref{98}--\sourcepageref{101}, \sourcepageref{255}--\sourcepageref{258}}
\idxsub{時間尺度（一般）}{\sourcepageref{94}}
\idxsub{時間尺度（量子Brown運動）}{\sourcepageref{197}}
\idxsub{時間尺度（散乱）}{\sourcepageref{127}, \sourcepageref{131}, \sourcepageref{134}}
\idxsub{仮想的}{\sourcepageref{98}--\sourcepageref{101}}
\idxsub{古典的雑音との比較}{\sourcepageref{95}--\sourcepageref{98}, \sourcepageref{302}--\sourcepageref{304}}
\idxsub{制御可能性との比較}{\sourcepageref{301}--\sourcepageref{302}}
\idxsub{散逸との比較}{\sourcepageref{93}--\sourcepageref{95}}
\idxentry{デコヒーレンス率}{\sourcepageref{70}, \sourcepageref{94}}
\idxsub{環境散乱におけるもの}{\sourcepageref{131}, \sourcepageref{134}}
\idxsub{微小管におけるもの}{\sourcepageref{371}--\sourcepageref{374}}
\idxsub{ニューロンにおけるもの}{\sourcepageref{368}--\sourcepageref{371}}
\idxsub{Caldeira–Leggett模型におけるもの}{\sourcepageref{193}}
\idxsub{スピン・ボソン模型におけるもの}{\sourcepageref{216}}
\idxsub{飽和}{\sourcepageref{132}, \sourcepageref{203}}
\idxentry{デコヒーレンス時間尺度}{$\to$ デコヒーレンス率}
\idxentry{デコヒーレンスのない部分空間}{\sourcepageref{79}--\sourcepageref{80}, \sourcepageref{321}--\sourcepageref{327}}
\idxsub{実験的実現}{\sourcepageref{325}--\sourcepageref{326}}
\idxsub{構造と大きさ}{\sourcepageref{322}--\sourcepageref{325}}
\idxentry{遅延選択実験}{\sourcepageref{100}}
\idxentry{位相関係の非局在化}{\sourcepageref{69}, \sourcepageref{98}--\sourcepageref{101}}
\idxentry{非局在モード}{\sourcepageref{175}}
\idxentry{密度行列}{\sourcepageref{33}--\sourcepageref{49}}
\idxsub{基底の曖昧さ}{\sourcepageref{41}--\sourcepageref{42}}
\idxsub{解釈}{\sourcepageref{36}, \sourcepageref{39}, \sourcepageref{48}--\sourcepageref{49}}
\idxsub{混合状態のもの}{\sourcepageref{36}--\sourcepageref{41}}
\idxsub{純粋状態のもの}{\sourcepageref{34}--\sourcepageref{36}}
\idxsub{純粋度}{\sourcepageref{40}}
\idxsub{縮約されたもの}{\sourcepageref{44}--\sourcepageref{49}}
\idxsub{von Neumannエントロピー}{\sourcepageref{41}}
\idxentry{密度演算子}{$\to$ 密度行列}
\idxentry{位相緩和}{\sourcepageref{97}, \sourcepageref{270}, \sourcepageref{282}}
\idxentry{微分断面積}{\sourcepageref{123}}
\idxentry{直接ピーク}{\sourcepageref{107}}
\idxentry{散逸}{\sourcepageref{93}--\sourcepageref{95}, \sourcepageref{188}--\sourcepageref{191}}
\idxentry{散逸核}{\sourcepageref{181}}
\idxentry{識別可能性}{\sourcepageref{32}--\sourcepageref{33}, \sourcepageref{47}--\sourcepageref{48}}
\idxentry{二重スリット実験}{\sourcepageref{23}, \sourcepageref{60}, \sourcepageref{258}--\sourcepageref{259}, \sourcepageref{355}}
\idxentry{二重井戸ポテンシャル}{\sourcepageref{174}, \sourcepageref{273}}
\idxentry{動力学的分離}{\sourcepageref{10}, \sourcepageref{326}--\sourcepageref{327}, \sourcepageref{375}}
\idxentry{動力学的対称性}{\sourcepageref{80}}
\idxentry{動力学的に分離された部分空間}{\sourcepageref{327}}
\idxentry{Dyson級数}{\sourcepageref{210}}
\idxentry{固有値・固有状態リンク}{\sourcepageref{59}, \sourcepageref{331}}
\idxentry{アインセレクション}{$\to$ 環境誘起超選択}
\idxentry{熱放射の放出}{\sourcepageref{115}, \sourcepageref{265}--\sourcepageref{267}}
\idxentry{集団}{}
\idxsub{平均}{\sourcepageref{43}, \sourcepageref{96}}
\idxsub{古典的（真正）}{\sourcepageref{20}}
\idxsub{位相緩和}{\sourcepageref{97}, \sourcepageref{270}, \sourcepageref{282}}
\idxsub{非真正}{\sourcepageref{49}}
\idxsub{集団解釈}{\sourcepageref{19}}
\idxsub{物理的}{\sourcepageref{43}--\sourcepageref{44}, \sourcepageref{96}, \sourcepageref{303}}
\idxsub{幅}{\sourcepageref{141}--\sourcepageref{147}}
\idxentry{もつれ}{\sourcepageref{4}, \sourcepageref{28}--\sourcepageref{33}}
\idxsub{識別可能性との関係}{\sourcepageref{32}--\sourcepageref{33}}
\idxsub{経路情報との関係}{\sourcepageref{60}--\sourcepageref{68}}
\idxsub{資源として}{\sourcepageref{7}}
\idxsub{定量化}{\sourcepageref{32}--\sourcepageref{33}}
\idxentry{エントロピー}{\sourcepageref{41}, \sourcepageref{313}}
\idxsub{Shannon}{\sourcepageref{86}}
\idxsub{熱力学的}{\sourcepageref{41}}
\idxsub{von Neumann}{\sourcepageref{41}}
\idxentry{エンバリアンス}{\sourcepageref{341}--\sourcepageref{343}}
\idxentry{証人としての環境}{\sourcepageref{86}}
\idxentry{環境設計}{\sourcepageref{326}--\sourcepageref{327}}
\idxentry{環境自己相関関数}{\sourcepageref{157}, \sourcepageref{181}}
\idxentry{環境援助不変性}{\sourcepageref{341}--\sourcepageref{343}}
\idxentry{環境誘起超選択}{\sourcepageref{6}, \sourcepageref{71}--\sourcepageref{85}}
\idxsub{解釈に対する含意}{\sourcepageref{329}--\sourcepageref{358}}
\idxsub{予測可能性のふるいとの関係}{\sourcepageref{82}--\sourcepageref{83}}
\idxsub{可換性基準}{\sourcepageref{77}}
\idxsub{一般の場合}{\sourcepageref{81}--\sourcepageref{83}}
\idxsub{Bohm力学におけるもの}{\sourcepageref{356}--\sourcepageref{357}}
\idxsub{様相解釈におけるもの}{\sourcepageref{345}--\sourcepageref{346}}
\idxsub{物理的収縮理論におけるもの}{\sourcepageref{349}--\sourcepageref{351}}
\idxsub{相対状態解釈におけるもの}{\sourcepageref{337}--\sourcepageref{344}}
\idxsub{Copenhagen解釈におけるもの}{\sourcepageref{334}--\sourcepageref{335}}
\idxsub{デコヒーレンスの量子極限におけるもの}{\sourcepageref{81}}
\idxsub{量子測定極限におけるもの}{\sourcepageref{76}--\sourcepageref{80}}
\idxsub{電荷のもの}{\sourcepageref{84}}
\idxsub{位置のもの}{\sourcepageref{83}}
\idxsub{安定性基準}{\sourcepageref{73}}
\idxentry{環境による監視}{\sourcepageref{65}--\sourcepageref{68}, \sourcepageref{78}}
\idxentry{環境散乱}{\sourcepageref{65}, \sourcepageref{115}--\sourcepageref{151}}
\idxentry{環境の自己相関}{\sourcepageref{156}, \sourcepageref{163}}
\idxentry{認識論的見方}{\sourcepageref{19}}
\idxentry{EPR}{\sourcepageref{6}, \sourcepageref{31}}
\idxentry{実存的解釈}{\sourcepageref{344}}
\idxentry{忠実な測定}{\sourcepageref{75}}
\idxentry{偽のデコヒーレンス}{\sourcepageref{96}}
\idxentry{耐故障性}{\sourcepageref{320}--\sourcepageref{321}}
\idxentry{第一の介入}{\sourcepageref{58}, \sourcepageref{363}}
\idxentry{磁束量子}{\sourcepageref{272}}
\idxentry{Fokker–Planck方程式}{\sourcepageref{184}}
\idxentry{自由意志}{\sourcepageref{366}}
\idxentry{ゲート操作}{\sourcepageref{296}}
\idxentry{Gauss型密度行列}{\sourcepageref{140}}
\idxentry{Gauss波束}{\sourcepageref{116}}
\idxentry{Groverの探索アルゴリズム}{\sourcepageref{300}}
\idxentry{GRW理論}{$\to$ 物理的収縮理論}
\idxentry{誘導方程式}{\sourcepageref{355}}
\idxentry{Hadamard回転}{\sourcepageref{317}}
\idxentry{Hamiltonian}{}
\idxsub{相互作用}{\sourcepageref{76}}
\idxsub{Lambシフトを含むもの}{\sourcepageref{155}}
\idxsub{開放系への分解}{\sourcepageref{76}}
\idxsub{自己Hamiltonian}{\sourcepageref{76}}
\idxentry{調和振動子}{\sourcepageref{117}, \sourcepageref{171}}
\idxentry{Heisenberg切断}{\sourcepageref{27}, \sourcepageref{335}}
\idxentry{発見的な擬制}{\sourcepageref{376}}
\idxentry{隠れた変数理論}{\sourcepageref{17}--\sourcepageref{18}, \sourcepageref{355}}
\idxentry{Hilbert空間}{\sourcepageref{15}}
\idxentry{Hyperion}{\sourcepageref{118}}
\idxentry{理想的量子測定}{\sourcepageref{52}}
\idxentry{無知で解釈できる混合}{\sourcepageref{42}}
\idxentry{非真正混合}{\sourcepageref{49}, \sourcepageref{69}, \sourcepageref{333}}
\idxentry{無限後退}{\sourcepageref{362}}
\idxentry{情報は物理的である}{\sourcepageref{294}, \sourcepageref{344}}
\idxentry{相互作用Hamiltonian}{\sourcepageref{76}}
\idxsub{対角分解}{\sourcepageref{78}}
\idxentry{相互作用描像}{\sourcepageref{379}--\sourcepageref{382}}
\idxentry{相互作用描像のLiouville–von Neumann方程式}{\sourcepageref{380}}
\idxentry{干渉}{\sourcepageref{47}--\sourcepageref{48}}
\idxsub{識別可能性との関係}{\sourcepageref{47}--\sourcepageref{48}}
\idxsub{実験}{\sourcepageref{23}--\sourcepageref{26}, \sourcepageref{243}--\sourcepageref{291}}
\idxsub{Bose–Einstein凝縮体におけるもの}{\sourcepageref{282}--\sourcepageref{283}}
\idxsub{$\mathrm C_{70}$回折におけるもの}{\sourcepageref{258}--\sourcepageref{270}}
\idxsub{SQUIDにおけるもの}{\sourcepageref{275}--\sourcepageref{279}}
\idxsub{局所的抑制}{\sourcepageref{68}--\sourcepageref{71}}
\idxsub{観測不可能性}{\sourcepageref{50}, \sourcepageref{55}--\sourcepageref{57}, \sourcepageref{333}}
\idxsub{模様}{\sourcepageref{23}, \sourcepageref{24}, \sourcepageref{55}, \sourcepageref{60}}
\idxsub{項}{\sourcepageref{34}}
\idxentry{内部環境}{\sourcepageref{176}}
\idxentry{解釈基底}{\sourcepageref{339}}
\idxentry{イオントラップ}{\sourcepageref{301}, \sourcepageref{303}, \sourcepageref{305}, \sourcepageref{326}}
\idxentry{不可逆性}{}
\idxsub{情報喪失との関係}{\sourcepageref{101}}
\idxsub{量子消去との関係}{\sourcepageref{100}}
\idxsub{収縮のもの}{\sourcepageref{63}}
\idxsub{デコヒーレンスのもの}{\sourcepageref{69}, \sourcepageref{93}, \sourcepageref{98}, \sourcepageref{304}}
\idxsub{測定のもの}{\sourcepageref{69}, \sourcepageref{100}}
\idxentry{孤立系}{\sourcepageref{1}--\sourcepageref{8}}
\idxentry{Itô方程式}{\sourcepageref{349}}
\idxentry{Josephson効果}{\sourcepageref{271}}
\idxentry{Josephson接合}{\sourcepageref{272}}
\idxentry{微小管中のキンク状励起}{\sourcepageref{374}}
\idxentry{Kochen–Specker定理}{\sourcepageref{18}}
\idxentry{Kraus演算子形式}{\sourcepageref{110}--\sourcepageref{112}}
\idxentry{Lambシフトを含むHamiltonian}{\sourcepageref{155}}
\idxentry{Lindbladマスター方程式}{\sourcepageref{165}--\sourcepageref{169}}
\idxentry{Lindblad演算子}{\sourcepageref{167}}
\idxentry{Liouville–von Neumann方程式}{\sourcepageref{139}}
\idxsub{相互作用描像}{\sourcepageref{380}}
\idxentry{局所理論}{\sourcepageref{4}}
\idxentry{状態の局所性}{\sourcepageref{4}}
\idxentry{局在化}{\sourcepageref{115}--\sourcepageref{151}}
\idxentry{局在モード}{\sourcepageref{175}}
\idxentry{長波長極限}{}
\idxsub{環境散乱におけるもの}{\sourcepageref{130}--\sourcepageref{132}}
\idxsub{量子Brown運動におけるもの}{\sourcepageref{204}}
\idxentry{Loschmidtエコー}{\sourcepageref{88}}
\idxentry{低温環境}{\sourcepageref{175}}
\idxentry{多心解釈}{\sourcepageref{337}, \sourcepageref{377}}
\idxentry{多世界解釈}{\sourcepageref{337}}
\idxentry{写像}{\sourcepageref{173}--\sourcepageref{178}}
\idxsub{振動子環境とスピン環境との関係}{\sourcepageref{176}--\sourcepageref{178}}
\idxsub{中心系のもの}{\sourcepageref{173}--\sourcepageref{174}}
\idxsub{振動子環境のもの}{\sourcepageref{175}}
\idxsub{位相空間粒子のもの}{\sourcepageref{173}}
\idxsub{スピン環境のもの}{\sourcepageref{175}--\sourcepageref{176}}
\idxsub{スピン環境から振動子環境へのもの}{\sourcepageref{228}--\sourcepageref{237}}
\idxsub{二準位系のもの}{\sourcepageref{173}--\sourcepageref{174}}
\idxentry{Markov近似}{\sourcepageref{156}, \sourcepageref{163}}
\idxentry{マスター方程式}{\sourcepageref{153}--\sourcepageref{170}}
\idxsub{Born–Markov形}{\sourcepageref{155}--\sourcepageref{165}}
\idxsub{一般形式}{\sourcepageref{154}--\sourcepageref{155}}
\idxsub{Caldeira–Leggett形}{\sourcepageref{191}--\sourcepageref{194}}
\idxsub{Lindblad形}{\sourcepageref{165}--\sourcepageref{169}}
\idxsub{非Markov型}{\sourcepageref{169}--\sourcepageref{170}}
\idxentry{測定問題}{\sourcepageref{49}--\sourcepageref{60}}
\idxsub{主観的解決}{\sourcepageref{375}--\sourcepageref{378}}
\idxentry{測定型相互作用}{\sourcepageref{49}, \sourcepageref{67}}
\idxentry{測定}{\sourcepageref{15}}
\idxsub{現実化}{\sourcepageref{19}}
\idxsub{事後選択との関係}{\sourcepageref{21}}
\idxsub{Born–Pauli解釈}{\sourcepageref{19}}
\idxsub{収縮}{\sourcepageref{15}}
\idxsub{文脈依存性}{\sourcepageref{18}}
\idxsub{干渉}{\sourcepageref{21}--\sourcepageref{26}}
\idxsub{不可逆なもの}{\sourcepageref{69}, \sourcepageref{100}}
\idxsub{局所的なもの}{\sourcepageref{44}}
\idxsub{非可換性}{\sourcepageref{15}}
\idxsub{非破壊}{\sourcepageref{52}, \sourcepageref{74}}
\idxsub{観測量との関係}{\sourcepageref{15}}
\idxsub{重ね合わせのもの}{\sourcepageref{21}--\sourcepageref{26}}
\idxsub{射影的なもの}{\sourcepageref{15}}
\idxsub{仮想的なもの}{\sourcepageref{101}}
\idxentry{記憶効果}{\sourcepageref{156}, \sourcepageref{163}, \sourcepageref{169}}
\idxentry{微小管}{\sourcepageref{367}, \sourcepageref{371}--\sourcepageref{374}}
\idxentry{混合状態}{\sourcepageref{19}, \sourcepageref{36}--\sourcepageref{41}}
\idxentry{混合}{}
\idxsub{無知で解釈できるもの}{\sourcepageref{42}}
\idxsub{非真正}{\sourcepageref{49}, \sourcepageref{69}, \sourcepageref{333}}
\idxsub{真正}{\sourcepageref{20}, \sourcepageref{42}}
\idxentry{様相解釈}{\sourcepageref{344}--\sourcepageref{347}}
\idxsub{環境誘起超選択との関係}{\sourcepageref{345}--\sourcepageref{346}}
\idxsub{Schmidt分解}{\sourcepageref{346}--\sourcepageref{347}}
\idxentry{模型}{}
\idxsub{標準的}{\sourcepageref{171}--\sourcepageref{241}}
\idxsub{量子Brown運動のもの}{\sourcepageref{178}--\sourcepageref{206}}
\idxsub{スピン・ボソン}{\sourcepageref{207}--\sourcepageref{222}}
\idxsub{スピン環境}{\sourcepageref{88}--\sourcepageref{93}, \sourcepageref{222}--\sourcepageref{237}}
\idxentry{相互情報量}{\sourcepageref{86}}
\idxentry{ミエリン}{\sourcepageref{368}}
\idxentry{ニューロン}{\sourcepageref{365}--\sourcepageref{371}, \sourcepageref{375}--\sourcepageref{378}}
\idxentry{NMR}{\sourcepageref{88}, \sourcepageref{97}, \sourcepageref{301}}
\idxentry{複製不可能定理}{\sourcepageref{16}, \sourcepageref{306}}
\idxentry{信号伝達不能定理}{\sourcepageref{31}}
\idxentry{雑音核}{\sourcepageref{181}}
\idxentry{非Hermite観測量}{\sourcepageref{53}}
\idxentry{非Markov動力学}{\sourcepageref{169}--\sourcepageref{170}}
\idxentry{観測量の非可換性}{\sourcepageref{15}, \sourcepageref{54}}
\idxentry{非破壊測定}{\sourcepageref{52}}
\idxsub{環境によるもの}{\sourcepageref{74}}
\idxentry{非局所性}{\sourcepageref{4}, \sourcepageref{6}, \sourcepageref{32}}
\idxentry{非ユニタリー動力学}{\sourcepageref{10}, \sourcepageref{28}, \sourcepageref{69}, \sourcepageref{111}, \sourcepageref{155}, \sourcepageref{290}, \sourcepageref{331}, \sourcepageref{348}--\sourcepageref{354}}
\idxentry{通常拡散係数}{\sourcepageref{185}, \sourcepageref{189}, \sourcepageref{191}, \sourcepageref{200}}
\idxentry{核磁気共鳴}{\sourcepageref{88}, \sourcepageref{97}, \sourcepageref{301}}
\idxentry{冗長性による客観化}{\sourcepageref{88}}
\idxentry{観測量}{\sourcepageref{15}}
\idxentry{観測}{}
\idxsub{量子力学におけるもの}{\sourcepageref{359}--\sourcepageref{361}}
\idxsub{間接的なもの}{\sourcepageref{85}}
\idxentry{非対角項}{\sourcepageref{34}}
\idxentry{量子系の開放性}{\sourcepageref{3}--\sourcepageref{4}, \sourcepageref{69}, \sourcepageref{74}}
\idxentry{演算子和形式}{\sourcepageref{110}--\sourcepageref{112}}
\idxentry{光量子計算}{\sourcepageref{301}}
\idxentry{正統的解釈}{\sourcepageref{330}--\sourcepageref{336}}
\idxentry{振動子環境}{\sourcepageref{175}}
\idxsub{弱結合極限}{\sourcepageref{176}}
\idxentry{結果}{\sourcepageref{16}}
\idxsub{隠れた変数理論におけるもの}{\sourcepageref{17}}
\idxsub{結果の問題}{\sourcepageref{50}, \sourcepageref{57}--\sourcepageref{60}, \sourcepageref{331}--\sourcepageref{333}, \sourcepageref{362}}
\idxentry{過完備な状態集合}{\sourcepageref{53}, \sourcepageref{75}}
\idxentry{部分トレース}{\sourcepageref{44}, \sourcepageref{333}}
\idxsub{デコヒーレンスとの関係}{\sourcepageref{45}}
\idxsub{解釈}{\sourcepageref{44}, \sourcepageref{333}}
\idxentry{位相量子ビット}{\sourcepageref{281}}
\idxentry{位相関係}{\sourcepageref{69}}
\idxsub{非局在化}{\sourcepageref{69}, \sourcepageref{98}--\sourcepageref{101}}
\idxentry{位相反転誤り}{\sourcepageref{309}}
\idxentry{物理的収縮理論}{\sourcepageref{348}--\sourcepageref{354}}
\idxsub{デコヒーレンスとの競合}{\sourcepageref{351}}
\idxsub{実験的検証}{\sourcepageref{353}--\sourcepageref{354}}
\idxsub{優先基底問題}{\sourcepageref{349}--\sourcepageref{351}}
\idxsub{裾の問題}{\sourcepageref{351}--\sourcepageref{352}}
\idxentry{物理的集団}{\sourcepageref{43}--\sourcepageref{44}, \sourcepageref{96}, \sourcepageref{303}}
\idxentry{物理的実在}{\sourcepageref{15}, \sourcepageref{360}, \sourcepageref{375}}
\idxentry{Poincaré再帰}{\sourcepageref{93}, \sourcepageref{213}}
\idxentry{装置の指針}{\sourcepageref{51}}
\idxentry{指針基底}{\sourcepageref{75}--\sourcepageref{85}}
\idxentry{指針観測量}{\sourcepageref{77}--\sourcepageref{85}}
\idxentry{指針状態}{\sourcepageref{74}--\sourcepageref{85}}
\idxsub{解釈に対する含意}{\sourcepageref{329}--\sourcepageref{358}}
\idxsub{予測可能性のふるいとの関係}{\sourcepageref{82}}
\idxsub{可換性基準}{\sourcepageref{77}}
\idxsub{一般的選択法}{\sourcepageref{81}--\sourcepageref{83}}
\idxsub{電荷空間におけるもの}{\sourcepageref{84}}
\idxsub{様相解釈におけるもの}{\sourcepageref{345}--\sourcepageref{346}}
\idxsub{位相空間におけるもの}{\sourcepageref{194}--\sourcepageref{203}}
\idxsub{位置空間におけるもの}{\sourcepageref{83}}
\idxsub{相対状態解釈におけるもの}{\sourcepageref{337}--\sourcepageref{344}}
\idxsub{Copenhagen解釈におけるもの}{\sourcepageref{334}--\sourcepageref{335}}
\idxsub{デコヒーレンスの量子極限におけるもの}{\sourcepageref{81}}
\idxsub{量子測定極限におけるもの}{\sourcepageref{76}--\sourcepageref{80}}
\idxsub{予測可能性}{\sourcepageref{82}}
\idxsub{冗長な符号化}{\sourcepageref{87}}
\idxsub{安定性基準}{\sourcepageref{73}}
\idxentry{指針部分空間}{$\to$ デコヒーレンスのない部分空間}
\idxentry{実証主義的見方}{\sourcepageref{67}, \sourcepageref{360}}
\idxentry{事後選択}{\sourcepageref{21}}
\idxentry{予測可能性}{\sourcepageref{82}, \sourcepageref{376}}
\idxentry{予測可能性のふるい}{\sourcepageref{82}--\sourcepageref{83}}
\idxentry{優先基底}{\sourcepageref{55}, \sourcepageref{75}--\sourcepageref{85}}
\idxentry{優先観測量}{\sourcepageref{55}, \sourcepageref{74}--\sourcepageref{85}}
\idxentry{優先状態}{$\to$ 指針状態}
\idxentry{優先基底問題}{\sourcepageref{50}, \sourcepageref{53}--\sourcepageref{55}, \sourcepageref{71}, \sourcepageref{83}--\sourcepageref{85}}
\idxsub{様相解釈におけるもの}{\sourcepageref{344}--\sourcepageref{347}}
\idxsub{物理的収縮理論におけるもの}{\sourcepageref{349}--\sourcepageref{351}}
\idxsub{相対状態解釈におけるもの}{\sourcepageref{337}--\sourcepageref{340}}
\idxsub{Copenhagen解釈におけるもの}{\sourcepageref{334}--\sourcepageref{335}}
\idxentry{前測定}{\sourcepageref{53}}
\idxentry{確率}{}
\idxsub{エンバリアンスとの関係}{\sourcepageref{341}--\sourcepageref{343}}
\idxsub{Born}{\sourcepageref{35}}
\idxsub{古典的}{\sourcepageref{37}}
\idxsub{認識論的}{\sourcepageref{17}}
\idxsub{混合状態におけるもの}{\sourcepageref{36}}
\idxsub{相対状態解釈におけるもの}{\sourcepageref{340}--\sourcepageref{343}}
\idxsub{量子的}{\sourcepageref{15}--\sourcepageref{18}, \sourcepageref{331}}
\idxentry{確率振幅}{\sourcepageref{18}}
\idxentry{結果の問題}{\sourcepageref{50}, \sourcepageref{57}--\sourcepageref{60}, \sourcepageref{331}--\sourcepageref{333}, \sourcepageref{362}}
\idxentry{干渉の観測不可能性の問題}{\sourcepageref{50}, \sourcepageref{55}--\sourcepageref{57}, \sourcepageref{69}, \sourcepageref{333}}
\idxentry{真正混合}{\sourcepageref{20}, \sourcepageref{42}}
\idxentry{プロトフィラメント}{\sourcepageref{371}}
\idxentry{心身並行論}{\sourcepageref{361}, \sourcepageref{376}--\sourcepageref{378}}
\idxentry{純粋状態}{\sourcepageref{19}}
\idxentry{環境の純粋化}{\sourcepageref{109}}
\idxentry{純粋度}{\sourcepageref{40}, \sourcepageref{147}, \sourcepageref{201}}
\idxentry{量子アルゴリズム}{\sourcepageref{300}}
\idxentry{量子ビット}{\sourcepageref{88}, \sourcepageref{270}--\sourcepageref{282}, \sourcepageref{295}--\sourcepageref{297}}
\idxentry{量子的な脳}{\sourcepageref{359}--\sourcepageref{378}}
\idxentry{量子Brown運動}{\sourcepageref{178}--\sourcepageref{206}}
\idxsub{振動子・スピン模型との関係}{\sourcepageref{234}--\sourcepageref{237}}
\idxsub{Born–Markovマスター方程式}{\sourcepageref{178}--\sourcepageref{182}}
\idxsub{デコヒーレンス率}{\sourcepageref{193}}
\idxsub{動力学}{\sourcepageref{194}--\sourcepageref{203}}
\idxsub{厳密なマスター方程式}{\sourcepageref{206}}
\idxsub{調和振動子のマスター方程式}{\sourcepageref{182}--\sourcepageref{187}}
\idxsub{高温極限}{\sourcepageref{191}--\sourcepageref{194}}
\idxsub{Wigner表示}{\sourcepageref{184}, \sourcepageref{195}--\sourcepageref{197}}
\idxsub{模型の限界}{\sourcepageref{203}--\sourcepageref{205}}
\idxsub{位相空間デコヒーレンス}{\sourcepageref{194}--\sourcepageref{203}}
\idxsub{前Markov形のマスター方程式}{\sourcepageref{186}}
\idxentry{量子コンピューター}{\sourcepageref{30}, \sourcepageref{293}--\sourcepageref{328}}
\idxsub{脳との関係}{\sourcepageref{365}--\sourcepageref{374}}
\idxsub{デコヒーレンスと古典的雑音の比較}{\sourcepageref{302}--\sourcepageref{304}}
\idxsub{デコヒーレンスと制御可能性の比較}{\sourcepageref{301}--\sourcepageref{302}}
\idxsub{誤り訂正}{\sourcepageref{304}--\sourcepageref{321}}
\idxsub{耐故障性}{\sourcepageref{320}--\sourcepageref{321}}
\idxsub{イオントラップ}{\sourcepageref{301}, \sourcepageref{303}, \sourcepageref{305}, \sourcepageref{326}}
\idxsub{NMR}{\sourcepageref{301}}
\idxsub{光学式}{\sourcepageref{301}}
\idxsub{物理的実現}{\sourcepageref{300}--\sourcepageref{301}}
\idxsub{能力}{\sourcepageref{294}--\sourcepageref{297}}
\idxsub{読出し}{\sourcepageref{297}--\sourcepageref{298}}
\idxsub{冗長な符号化}{\sourcepageref{305}--\sourcepageref{307}, \sourcepageref{315}--\sourcepageref{320}}
\idxsub{記憶容量}{\sourcepageref{296}}
\idxentry{量子計算}{$\to$ 量子コンピューター}
\idxentry{量子相関}{\sourcepageref{4}, \sourcepageref{29}}
\idxentry{量子暗号}{\sourcepageref{30}}
\idxentry{量子Darwin主義}{\sourcepageref{86}}
\idxentry{量子ディスコード}{\sourcepageref{87}}
\idxentry{量子消去}{\sourcepageref{98}--\sourcepageref{101}, \sourcepageref{304}, \sourcepageref{316}}
\idxentry{量子誤り訂正}{\sourcepageref{304}--\sourcepageref{321}}
\idxsub{古典的誤り訂正との比較}{\sourcepageref{305}--\sourcepageref{307}}
\idxsub{誤りの離散化}{\sourcepageref{307}--\sourcepageref{311}}
\idxsub{耐故障型}{\sourcepageref{320}--\sourcepageref{321}}
\idxsub{3ビット符号}{\sourcepageref{315}--\sourcepageref{320}}
\idxentry{量子的全体論}{\sourcepageref{6}, \sourcepageref{33}, \sourcepageref{297}}
\idxentry{量子情報理論}{\sourcepageref{293}--\sourcepageref{328}}
\idxentry{デコヒーレンスの量子極限}{\sourcepageref{81}, \sourcepageref{228}}
\idxentry{量子Monte Carlo}{\sourcepageref{297}}
\idxentry{量子的観測者}{\sourcepageref{361}--\sourcepageref{365}}
\idxentry{量子並列性}{\sourcepageref{297}}
\idxentry{量子状態}{}
\idxsub{古典状態との比較}{\sourcepageref{14}--\sourcepageref{16}}
\idxsub{完全性}{\sourcepageref{16}}
\idxsub{概念と解釈}{\sourcepageref{14}--\sourcepageref{20}}
\idxsub{壊れやすさ}{\sourcepageref{16}, \sourcepageref{74}}
\idxsub{全体の状態}{\sourcepageref{29}}
\idxsub{Bell形}{\sourcepageref{29}}
\idxsub{混合}{\sourcepageref{19}}
\idxsub{存在論的地位}{\sourcepageref{18}--\sourcepageref{20}}
\idxsub{準備手順}{\sourcepageref{36}}
\idxsub{確率的性格}{\sourcepageref{16}--\sourcepageref{18}}
\idxsub{純粋}{\sourcepageref{19}}
\idxsub{実在論的解釈}{\sourcepageref{19}, \sourcepageref{348}, \sourcepageref{376}}
\idxsub{相対的}{\sourcepageref{29}}
\idxentry{量子電気機械系}{\sourcepageref{234}, \sourcepageref{284}--\sourcepageref{288}, \sourcepageref{336}}
\idxsub{デコヒーレンス}{\sourcepageref{287}--\sourcepageref{288}}
\idxsub{重ね合わせ状態}{\sourcepageref{285}--\sourcepageref{287}}
\idxentry{量子ジャンプ法}{\sourcepageref{168}}
\idxentry{量子測定極限}{\sourcepageref{76}--\sourcepageref{80}, \sourcepageref{228}}
\idxentry{量子から古典への移行}{\sourcepageref{2}, \sourcepageref{49}--\sourcepageref{60}}
\idxentry{量子軌道法}{\sourcepageref{168}}
\idxentry{量子ビット（qubit）}{\sourcepageref{88}, \sourcepageref{270}--\sourcepageref{282}, \sourcepageref{295}--\sourcepageref{297}}
\idxentry{Rabi周波数}{\sourcepageref{246}}
\idxentry{Rabi振動}{\sourcepageref{25}, \sourcepageref{246}--\sourcepageref{247}, \sourcepageref{276}}
\idxentry{Ramsey干渉法}{\sourcepageref{25}--\sourcepageref{26}, \sourcepageref{279}}
\idxentry{再コヒーレンス}{\sourcepageref{99}, \sourcepageref{255}--\sourcepageref{258}}
\idxentry{反跳するスリット}{\sourcepageref{61}}
\idxentry{コヒーレンスの再帰}{\sourcepageref{92}, \sourcepageref{213}}
\idxentry{Redfield方程式}{\sourcepageref{163}, \sourcepageref{187}}
\idxentry{縮約密度行列}{\sourcepageref{44}--\sourcepageref{49}, \sourcepageref{64}, \sourcepageref{68}}
\idxsub{デコヒーレンスとの関係}{\sourcepageref{45}}
\idxsub{導出}{\sourcepageref{45}--\sourcepageref{46}}
\idxsub{解釈}{\sourcepageref{48}--\sourcepageref{49}, \sourcepageref{333}}
\idxsub{正値性}{\sourcepageref{165}}
\idxsub{集団との比較}{\sourcepageref{48}--\sourcepageref{49}}
\idxentry{冗長性}{\sourcepageref{85}--\sourcepageref{88}, \sourcepageref{305}, \sourcepageref{336}, \sourcepageref{344}, \sourcepageref{378}}
\idxentry{情報の冗長な符号化}{}
\idxsub{古典コンピューターにおけるもの}{\sourcepageref{305}--\sourcepageref{307}}
\idxsub{量子コンピューターにおけるもの}{\sourcepageref{315}--\sourcepageref{320}}
\idxsub{環境におけるもの}{\sourcepageref{85}--\sourcepageref{88}}
\idxentry{相対状態}{\sourcepageref{29}, \sourcepageref{47}}
\idxentry{相対状態解釈}{\sourcepageref{336}--\sourcepageref{344}}
\idxsub{優先基底問題}{\sourcepageref{337}--\sourcepageref{340}}
\idxsub{確率}{\sourcepageref{340}--\sourcepageref{343}}
\idxsub{Schmidt分解}{\sourcepageref{339}}
\idxentry{緩和時間尺度}{\sourcepageref{93}}
\idxentry{リセット}{\sourcepageref{101}}
\idxentry{部分系への分割}{\sourcepageref{101}--\sourcepageref{103}}
\idxentry{共鳴$\pi/2$パルス}{\sourcepageref{247}}
\idxentry{ニューロンの静止・発火}{\sourcepageref{365}--\sourcepageref{371}}
\idxentry{可逆なデコヒーレンス}{\sourcepageref{98}--\sourcepageref{101}, \sourcepageref{255}--\sourcepageref{258}}
\idxentry{コヒーレンスの復活}{\sourcepageref{98}--\sourcepageref{101}, \sourcepageref{255}--\sourcepageref{258}}
\idxentry{rf-SQUID}{\sourcepageref{272}}
\idxentry{S行列}{\sourcepageref{119}}
\idxentry{散乱}{\sourcepageref{7}, \sourcepageref{45}, \sourcepageref{65}, \sourcepageref{115}--\sourcepageref{151}}
\idxsub{マスター方程式}{\sourcepageref{140}}
\idxsub{空気分子のもの}{\sourcepageref{136}--\sourcepageref{138}}
\idxsub{光子のもの}{\sourcepageref{132}--\sourcepageref{136}}
\idxentry{散乱定数}{\sourcepageref{131}}
\idxentry{Schmidt基底}{$\to$ Schmidt分解}
\idxentry{Schmidt分解}{\sourcepageref{104}--\sourcepageref{106}}
\idxsub{古典性との関係}{\sourcepageref{105}--\sourcepageref{106}, \sourcepageref{339}}
\idxsub{デコヒーレンスとの関係}{\sourcepageref{105}--\sourcepageref{106}}
\idxsub{指針状態との関係}{\sourcepageref{105}--\sourcepageref{106}}
\idxsub{縮退}{\sourcepageref{105}, \sourcepageref{347}}
\idxsub{密度行列に対するもの}{\sourcepageref{104}}
\idxsub{様相解釈におけるもの}{\sourcepageref{346}--\sourcepageref{347}}
\idxsub{相対状態解釈におけるもの}{\sourcepageref{339}}
\idxsub{一意性}{\sourcepageref{104}}
\idxentry{Schmidt状態}{$\to$ Schmidt分解}
\idxentry{Schrödinger方程式}{\sourcepageref{21}}
\idxsub{追加項を含むもの}{\sourcepageref{349}}
\idxsub{確率的}{\sourcepageref{169}, \sourcepageref{348}}
\idxentry{Schrödingerの猫}{\sourcepageref{2}--\sourcepageref{3}, \sourcepageref{57}--\sourcepageref{58}, \sourcepageref{364}}
\idxentry{第二の介入}{\sourcepageref{363}}
\idxentry{準古典的性質の選択}{\sourcepageref{83}--\sourcepageref{85}}
\idxentry{自己Hamiltonian}{\sourcepageref{76}}
\idxentry{Shannonエントロピー}{\sourcepageref{86}}
\idxentry{Shorの因数分解アルゴリズム}{\sourcepageref{300}}
\idxentry{短波長極限}{\sourcepageref{128}--\sourcepageref{129}}
\idxentry{「黙って計算せよ」解釈}{\sourcepageref{329}}
\idxentry{Solvay会議}{\sourcepageref{61}, \sourcepageref{329}}
\idxentry{空間的局在化}{\sourcepageref{74}, \sourcepageref{83}, \sourcepageref{115}--\sourcepageref{151}, \sourcepageref{265}, \sourcepageref{348}--\sourcepageref{351}}
\idxentry{特殊相対性理論}{\sourcepageref{31}}
\idxentry{スペクトル密度}{\sourcepageref{181}}
\idxsub{カットオフ}{\sourcepageref{188}}
\idxsub{有効なもの}{\sourcepageref{233}}
\idxsub{Lorentz–Drude形}{\sourcepageref{188}}
\idxsub{オーム型}{\sourcepageref{188}}
\idxentry{スピンエコー}{\sourcepageref{98}}
\idxentry{スピン環境}{\sourcepageref{175}--\sourcepageref{176}, \sourcepageref{222}--\sourcepageref{237}}
\idxentry{スピン・ボソン模型}{\sourcepageref{207}--\sourcepageref{222}}
\idxsub{スピン環境模型との関係}{\sourcepageref{233}--\sourcepageref{234}}
\idxsub{Born–Markovマスター方程式}{\sourcepageref{218}--\sourcepageref{222}}
\idxsub{簡略版}{\sourcepageref{208}--\sourcepageref{218}}
\idxentry{スピン環境模型}{\sourcepageref{222}--\sourcepageref{237}}
\idxsub{デコヒーレンスの量子極限}{\sourcepageref{228}}
\idxsub{量子測定極限}{\sourcepageref{228}}
\idxsub{弱結合極限}{\sourcepageref{228}--\sourcepageref{237}}
\idxsub{非Markov型}{\sourcepageref{237}}
\idxsub{単純な例}{\sourcepageref{223}--\sourcepageref{228}}
\idxentry{分枝の分裂}{\sourcepageref{337}, \sourcepageref{375}--\sourcepageref{378}}
\idxentry{自発的局在化模型}{\sourcepageref{348}--\sourcepageref{354}}
\idxentry{不気味な遠隔作用}{\sourcepageref{31}}
\idxentry{SQUID}{\sourcepageref{176}, \sourcepageref{270}--\sourcepageref{282}}
\idxsub{基本原理}{\sourcepageref{272}--\sourcepageref{275}}
\idxsub{デコヒーレンスの観測}{\sourcepageref{279}--\sourcepageref{282}}
\idxsub{物理的収縮理論との関係}{\sourcepageref{350}, \sourcepageref{354}}
\idxsub{Rabi振動}{\sourcepageref{276}}
\idxsub{重ね合わせ状態}{\sourcepageref{275}--\sourcepageref{279}}
\idxentry{安定性基準}{\sourcepageref{73}}
\idxentry{標準解釈}{\sourcepageref{330}--\sourcepageref{336}}
\idxentry{Stern–Gerlach実験}{\sourcepageref{21}--\sourcepageref{23}, \sourcepageref{37}}
\idxsub{可逆なもの}{\sourcepageref{99}}
\idxentry{確率的動力学的縮約}{\sourcepageref{348}--\sourcepageref{354}}
\idxentry{主観性}{\sourcepageref{360}}
\idxentry{部分系}{\sourcepageref{28}, \sourcepageref{101}--\sourcepageref{103}}
\idxentry{超伝導量子ビット}{\sourcepageref{270}--\sourcepageref{282}}
\idxentry{超伝導}{\sourcepageref{271}--\sourcepageref{272}}
\idxentry{超電流}{\sourcepageref{271}--\sourcepageref{272}}
\idxentry{重ね合わせ原理}{\sourcepageref{20}--\sourcepageref{28}}
\idxsub{解釈}{\sourcepageref{20}--\sourcepageref{21}}
\idxsub{適用範囲}{\sourcepageref{26}--\sourcepageref{28}}
\idxentry{重ね合わせ}{\sourcepageref{20}--\sourcepageref{28}}
\idxsub{増幅}{\sourcepageref{52}}
\idxsub{干渉実験との関係}{\sourcepageref{23}--\sourcepageref{26}}
\idxsub{コヒーレントなもの}{\sourcepageref{20}}
\idxsub{直接測定}{\sourcepageref{21}--\sourcepageref{23}}
\idxsub{実験的検証}{\sourcepageref{21}--\sourcepageref{26}}
\idxsub{Gauss状態のもの}{\sourcepageref{194}}
\idxsub{放射場のもの}{\sourcepageref{245}}
\idxsub{超電流のもの}{\sourcepageref{275}--\sourcepageref{279}}
\idxsub{重ね合わせにある量子ビット状態}{\sourcepageref{295}}
\idxsub{再局在化}{\sourcepageref{99}, \sourcepageref{255}--\sourcepageref{258}}
\idxsub{古典的集団との比較}{\sourcepageref{20}}
\idxentry{超選択則}{\sourcepageref{6}, \sourcepageref{73}}
\idxentry{環境の対称化}{\sourcepageref{326}}
\idxentry{系・観測者二元論}{\sourcepageref{336}}
\idxentry{T行列}{\sourcepageref{122}, \sourcepageref{123}}
\idxentry{裾の問題}{\sourcepageref{351}--\sourcepageref{352}}
\idxentry{Talbot–Lau効果}{\sourcepageref{260}--\sourcepageref{261}}
\idxentry{熱浴}{\sourcepageref{171}}
\idxentry{3ビット符号}{\sourcepageref{315}--\sourcepageref{320}}
\idxentry{閾値定理}{\sourcepageref{321}}
\idxentry{時間尺度}{}
\idxsub{コヒーレントな広がり}{\sourcepageref{117}}
\idxsub{デコヒーレンス（一般）}{\sourcepageref{70}, \sourcepageref{94}}
\idxsub{量子Brown運動のデコヒーレンス}{\sourcepageref{197}}
\idxsub{Caldeira–Leggett模型のデコヒーレンス}{\sourcepageref{185}}
\idxsub{スピン・ボソン模型のデコヒーレンス}{\sourcepageref{216}}
\idxsub{コヒーレンスの再帰}{\sourcepageref{93}}
\idxsub{緩和}{\sourcepageref{93}}
\idxsub{散乱誘起デコヒーレンス}{\sourcepageref{127}, \sourcepageref{131}, \sourcepageref{134}}
\idxentry{トレース操作}{\sourcepageref{34}--\sourcepageref{36}}
\idxsub{デコヒーレンスとの関係}{\sourcepageref{45}}
\idxsub{解釈}{\sourcepageref{36}, \sourcepageref{44}, \sourcepageref{333}}
\idxsub{部分トレース}{\sourcepageref{44}, \sourcepageref{333}}
\idxentry{トレース規則}{\sourcepageref{36}, \sourcepageref{46}}
\idxsub{混合状態に対するもの}{\sourcepageref{38}}
\idxentry{チューブリン}{\sourcepageref{371}}
\idxentry{トンネル}{\sourcepageref{174}, \sourcepageref{271}, \sourcepageref{277}}
\idxentry{トンネル行列要素}{\sourcepageref{174}}
\idxentry{不確定性原理}{\sourcepageref{62}, \sourcepageref{202}, \sourcepageref{366}}
\idxentry{仮想的デコヒーレンス}{\sourcepageref{98}--\sourcepageref{101}}
\idxentry{仮想的測定}{\sourcepageref{101}}
\idxentry{von Neumann連鎖}{\sourcepageref{361}--\sourcepageref{365}}
\idxentry{von Neumannエントロピー}{\sourcepageref{41}}
\idxentry{von Neumann測定図式}{\sourcepageref{50}--\sourcepageref{53}, \sourcepageref{361}--\sourcepageref{362}}
\idxentry{波束}{\sourcepageref{18}}
\idxsub{コヒーレントな広がり}{\sourcepageref{116}}
\idxsub{Gauss型}{\sourcepageref{116}}
\idxentry{弱結合極限}{\sourcepageref{176}}
\idxentry{量子誤り訂正における重み}{\sourcepageref{310}}
\idxentry{経路情報}{\sourcepageref{60}--\sourcepageref{68}, \sourcepageref{263}--\sourcepageref{265}}
\idxsub{もつれとの関係}{\sourcepageref{60}--\sourcepageref{68}}
\idxsub{環境による監視との関係}{\sourcepageref{65}--\sourcepageref{68}}
\idxentry{Wigner表示}{\sourcepageref{106}--\sourcepageref{109}}
\idxsub{直接ピーク}{\sourcepageref{107}}
\idxsub{量子Brown運動におけるもの}{\sourcepageref{184}, \sourcepageref{195}--\sourcepageref{197}}
\idxsub{振動模様}{\sourcepageref{107}}
\idxentry{Wignerの友人}{\sourcepageref{364}--\sourcepageref{365}}
\end{multicols}
\endgroup
